\documentclass[11pt,a4paper]{article}

\usepackage[utf8]{inputenc}
\usepackage[T1]{fontenc}
\usepackage[margin=2.5cm]{geometry}
\usepackage{booktabs}
\usepackage{array}
\usepackage{xcolor}
\usepackage{hyperref}
\usepackage{amsmath}
\usepackage{amssymb}
\usepackage{enumitem}
\usepackage{multirow}
\usepackage{caption}
\usepackage{parskip}

\definecolor{darkblue}{RGB}{0,50,120}
\hypersetup{colorlinks=true,linkcolor=darkblue,urlcolor=darkblue}

\newcolumntype{L}[1]{>{\raggedright\arraybackslash}p{#1}}
\newcolumntype{C}[1]{>{\centering\arraybackslash}p{#1}}
\newcolumntype{R}[1]{>{\raggedleft\arraybackslash}p{#1}}

\title{%
  \textbf{Betweenness Centrality Prediction with Hyper-Connected GNNs}\\[0.4em]
  \large Project Progress Report
}
\author{Research Log Synthesis}
\date{March 11, 2026}

\begin{document}
\maketitle
\tableofcontents
\newpage

%------------------------------------------------------------
\section{Research Question}
%------------------------------------------------------------

Can deep GNNs predict betweenness centrality (BC) inductively,
and do HC / mHC / mHC-lite improve depth scaling and
out-of-distribution (OOD) ranking quality
compared with strong deep-GNN baselines?

\textbf{Primary evaluation metrics:}
\begin{itemize}[noitemsep]
  \item Ranking: Kendall $\tau$, Spearman $\rho$
  \item Top-K: Precision@K, NDCG@K
\end{itemize}

\textbf{Protocol (inductive):} train on small synthetic graphs
(100 nodes, ER/BA/SBM), evaluate on larger unseen graphs
(500 nodes, same families = ID; WS/RGG families = OOD).

%------------------------------------------------------------
\section{Stage 1 --- Full 5-Seed Depth Sweep (2026-03-06)}
%------------------------------------------------------------

\subsection{Setup}
\begin{table}[h]
\centering
\begin{tabular}{L{4.5cm} L{7.5cm}}
\toprule
Seeds & 5 (0--4) \\
Models & 19 variants: GCN, SAGE, GAT, GIN, GCNII, APPNP, JKNet,
         HC, mHC, mHC-lite \\
Depths & $L \in \{2,4,8,16,32\}$ \\
Train graphs & 160 (100 nodes, ER/BA/SBM) \\
Val / Test ID / OOD & 30 / 30 / 30 (500 nodes) \\
Budget & 50 epochs, patience 10, LR $10^{-3}$, WD $10^{-5}$ \\
\bottomrule
\end{tabular}
\end{table}

\subsection{Best points from the frozen sweep}
\begin{table}[h]
\centering
\begin{tabular}{L{4.5cm} L{2.5cm} C{1.2cm} R{3.5cm}}
\toprule
Category & Model & $L$ & OOD Kendall \\
\midrule
Best overall OOD     & \texttt{hc\_gcnii}  & 2 & $0.3824 \pm 0.0253$ \\
Best mHC point       & \texttt{mhc\_gcn}   & 2 & $0.3812 \pm 0.0238$ \\
Strongest GCN gain   & \texttt{hc\_gcn}    & 4 & $0.3809 \pm 0.0386$ \\
Best deep baseline   & \texttt{appnp}      & 32 & $0.3544 \pm 0.0258$ \\
\bottomrule
\end{tabular}
\end{table}

\subsection{Conclusions}
HC / mHC improve OOD consistently; gains concentrate at $L \leq 8$.
No ultra-deep scaling claim is supported.
Two families retained: \textbf{GCN} (clearest mHC signal)
and \textbf{GCNII} (deep-backbone stress test).

%------------------------------------------------------------
\section{Stage 2 --- Targeted Hyperparameter Search (2026-03-06)}
%------------------------------------------------------------

Optuna proxy search (15 trials, validation-only objective):

\begin{itemize}
  \item \textbf{GCN Trial 3} (obj. 0.8605): LR $2.26\times10^{-4}$,
        WD $2.10\times10^{-5}$, dropout 0.1, $n_\text{str}=2$,
        $\tau=0.2$, iters 20.
  \item \textbf{GCNII Trial 11} (obj. 0.8564): LR $1.07\times10^{-3}$,
        WD $4.34\times10^{-5}$, $n_\text{str}=4$, $\tau=0.1$,
        GCNII $\alpha=0.1$, $\theta=0.5$.
\end{itemize}

GCNII failure at L16 persists after tuning.
GCN + mHC is the strongest candidate.

%------------------------------------------------------------
\section{Stage 3 --- Multi-Seed Confirmation (2026-03-07)}
%------------------------------------------------------------

90 runs per family (3 candidate trials $\times$ 3 depths $\times$ 5 seeds).

\textbf{GCN:} mHC-GCN consistently gains on OOD Kendall at L4/L8
(strongest: $+0.071$ at t3/L4). Single negative case: t13/L2 ($-0.006$).

\textbf{GCNII:} mHC helps at L2/L8 but degrades at L16
(deltas: $-0.039$, $-0.012$, $-0.047$).

\textit{Decision: the GCN family becomes the primary line of analysis from this point; other backbones remain secondary reference points.}

%------------------------------------------------------------
\section{Stage 4 --- GCN Ablation Study (2026-03-07)}
%------------------------------------------------------------

\subsection{Results at L8 (tuned GCN regime, 5 seeds)}
\begin{table}[h]
\centering
\begin{tabular}{L{3.5cm} L{3.5cm} R{2.5cm} R{2.5cm}}
\toprule
Block & Variant & Val Kendall & OOD Kendall \\
\midrule
Structure & gcn\_ref          & 0.8332 & 0.3574 \\
Structure & mhc\_gcn\_ref     & 0.8909 & 0.3888 \\
Structure & \textbf{hc\_gcn\_ref} & \textbf{0.9174} & \textbf{0.4054} \\
Structure & dynamic only      & 0.8882 & 0.3904 \\
Structure & static only       & 0.8832 & 0.3855 \\
Structure & $n=4$ streams     & 0.8809 & 0.3791 \\
\midrule
Sinkhorn  & $\tau=0.05$       & 0.8873 & 0.3923 \\
Sinkhorn  & $\tau=0.10$       & 0.8932 & 0.3973 \\
Sinkhorn  & $\tau=0.20$       & 0.8952 & 0.3911 \\
Sinkhorn  & iters $=10$       & 0.8886 & 0.3906 \\
\bottomrule
\end{tabular}
\end{table}

\subsection{Interpretation}
HC-GCN dominates. No mHC variant exceeds HC regardless of $\tau$ tuning.
$n=2$ streams is the best setting in this study. Dynamic-only and
static-only variants both remain functional, while the hybrid setting
is the most stable default. The useful mechanism appears to be flexible
multi-stream routing, more than the doubly-stochastic Sinkhorn
constraint itself.

%------------------------------------------------------------
\section{Stage 5 --- Final GCN Family Comparison (2026-03-09)}
%------------------------------------------------------------

\begin{table}[h]
\centering
\caption{Final GCN family --- test OOD results (5 seeds)}
\begin{tabular}{L{3cm} C{0.8cm} R{2.2cm} R{2.2cm} R{2.2cm} R{2.2cm}}
\toprule
Variant & $L$ & Kendall & Spearman & P@10 & NDCG@10 \\
\midrule
gcn            & 4 & $0.302\pm0.045$ & $0.423\pm0.063$ & 0.247 & 0.513 \\
gcn            & 8 & $0.360\pm0.021$ & $0.506\pm0.028$ & 0.264 & 0.563 \\
\midrule
hc\_gcn        & 4 & $0.393\pm0.031$ & $0.543\pm0.041$ & $\mathbf{0.291}$ & $\mathbf{0.571}$ \\
hc\_gcn        & 8 & $\mathbf{0.406\pm0.041}$ & $\mathbf{0.558\pm0.054}$ & 0.287 & 0.567 \\
\midrule
mhc\_gcn       & 4 & $0.374\pm0.043$ & $0.521\pm0.055$ & 0.274 & 0.566 \\
mhc\_gcn       & 8 & $0.396\pm0.038$ & $0.548\pm0.049$ & 0.278 & $\mathbf{0.568}$ \\
\midrule
mhc\_lite\_gcn & 4 & $0.376\pm0.042$ & $0.523\pm0.055$ & 0.271 & 0.553 \\
mhc\_lite\_gcn & 8 & $0.395\pm0.039$ & $0.544\pm0.050$ & 0.282 & 0.561 \\
\bottomrule
\end{tabular}
\end{table}

\textbf{Key findings:}
\begin{enumerate}
  \item \textbf{Val:} HC gains are clearly learned in validation and remain visible on
        Kendall, Precision@10, and NDCG@10. At $L=8$, HC-GCN reaches
        $0.918$ Kendall, $0.919$ Precision@10, and $0.995$ NDCG@10,
        versus $0.833$, $0.851$, and $0.976$ for plain GCN. Depth has little
        effect on plain GCN, while mHC and mHC-lite improve from $L=4$ to $L=8$.
  \item \textbf{Test ID:} All models remain in a narrow band on in-distribution test
        Kendall (roughly $0.82$--$0.83$), suggesting that hyper-connections bring
        little benefit on the known distribution.
  \item \textbf{Test OOD:} Gains are substantial on out-of-distribution graphs:
        HC-GCN $>$ mHC-GCN $\approx$ mHC-lite-GCN $>$ GCN on the main ranking
        metrics. On top-K at $10$, HC is strongest on Precision@10, while mHC is
        marginally strongest on NDCG@10.
\end{enumerate}

%------------------------------------------------------------
\section{Stage 6 --- Residual Baseline (2026-03-10)}
%------------------------------------------------------------

\begin{table}[h]
\centering
\caption{Residual baseline --- OOD Kendall (5 seeds)}
\begin{tabular}{L{3.5cm} C{1.2cm} R{3.5cm} R{3cm}}
\toprule
Variant & $L$ & OOD Kendall & Relative gain \\
\midrule
gcn             & 4 & $0.305\pm0.046$ & --- \\
gcn\_residual   & 4 & $0.352\pm0.033$ & $+15.2\%$ \\
hc\_gcn         & 4 & $0.395\pm0.037$ & $+29.3\%$ \\
mhc\_gcn        & 4 & $0.380\pm0.049$ & $+24.4\%$ \\
\midrule
gcn             & 8 & $0.359\pm0.024$ & --- \\
gcn\_residual   & 8 & $0.369\pm0.035$ & $+2.6\%$ \\
hc\_gcn         & 8 & $\mathbf{0.411\pm0.040}$ & $+14.5\%$ \\
mhc\_gcn        & 8 & $0.389\pm0.046$ & $+8.4\%$ \\
\bottomrule
\end{tabular}
\end{table}

At L8: residual gives only $+2.6\%$; HC-GCN gives $+14.5\%$.
The hyper-connection gain is architectural, not reducible to
skip connections.

%------------------------------------------------------------
\section{Stage 7 --- Zero-Shot Real-Graph Transfer (2026-03-09/10)}
%------------------------------------------------------------

\subsection{Setup}
Models trained only on synthetic graphs evaluated zero-shot on:
\begin{itemize}[noitemsep]
  \item \textbf{Cora} (2708 nodes, 5278 edges, $\bar{d}\approx3.9$)
  \item \textbf{CiteSeer} (3327 nodes, 4552 edges, $\bar{d}\approx2.7$)
  \item \textbf{PubMed} (19717 nodes, 44324 edges, $\bar{d}\approx4.5$)
\end{itemize}

\subsection{Results --- Kendall $\tau$ (mean $\pm$ std, 5 seeds)}
\begin{table}[h]
\centering
\begin{tabular}{L{3.8cm} R{3cm} R{3cm} R{3cm}}
\toprule
Variant & Cora & CiteSeer & PubMed \\
\midrule
gcn\_l8            & $0.350\pm0.053$ & $0.067\pm0.118$ & $0.539\pm0.025$ \\
gcn\_residual\_l8  & $0.355\pm0.062$ & $0.069\pm0.080$ & $0.605\pm0.042$ \\
hc\_gcn\_l8        & $0.366\pm0.034$ & $0.056\pm0.029$ & $0.595\pm0.028$ \\
mhc\_gcn\_l8       & $\mathbf{0.405\pm0.036}$ & $\mathbf{0.151\pm0.083}$ & $0.598\pm0.030$ \\
mhc\_lite\_gcn\_l8 & $0.380\pm0.037$ & $0.125\pm0.085$ & $\mathbf{0.611\pm0.030}$ \\
\bottomrule
\end{tabular}
\end{table}

\subsection{Per-graph interpretation}

\textbf{Cora} ($\bar{d}\approx3.9$). All hyper-connected variants improve.
BC is not severely skewed. mHC-GCN leads on all metrics
($+15.5\%$ Kendall, $+43.8\%$ P@top5\%).
HC also gains ($+4.5\%$) without regression.

\textbf{CiteSeer} ($\bar{d}\approx2.7$ --- extreme sparsity).
Most informative graph. BC concentrates on a small number of bridge nodes.
HC-GCN \emph{regresses} $-16\%$ in Kendall vs.\ plain GCN,
while achieving best NDCG@10 ($+34\%$):
HC learns to route information toward bridges (correct top-K)
but destroys global ranking of other nodes.
The Sinkhorn constraint of mHC prevents this collapse:
$+127\%$ Kendall gain.
The routing matrix analysis (Section~\ref{sec:matrix-analysis}) provides
a mechanistic account of this divergence.

\textbf{PubMed} (19717 nodes).
Gains more moderate. GCN+residual is very competitive ($+12.2\%$).
mHC-lite leads ($+13.3\%$): full Sinkhorn overhead not worthwhile
at this scale; lighter variant offers better trade-off.
Precision@10 saturates at 0.5 for all models (non-discriminative).

\subsection{Routing collapse on sparse BC}

The HC regression on CiteSeer reveals a regime-dependent property:
\begin{center}
\begin{tabular}{L{5.5cm} L{5.5cm}}
\toprule
\textbf{HC (unconstrained)} & \textbf{mHC (Sinkhorn)} \\
\midrule
Best on synthetic BC (moderate skew) & Best on real citation graphs \\
High routing freedom & Proportional stream balancing \\
Collapses on sparse/skewed BC & No regression on any tested graph \\
\bottomrule
\end{tabular}
\end{center}

%------------------------------------------------------------
\section{Stage 8 --- Routing Matrix Analysis (2026-03-11)}
\label{sec:matrix-analysis}
%------------------------------------------------------------

\subsection{Motivation}

Performance alone tells us \emph{which} model wins, but not \emph{why}.
To explain the synthetic vs.\ real-graph ranking inversion
(HC best on synthetic, mHC/mHC-lite best on citation graphs),
we inspected the learned inter-stream residual routing matrices
$H_{\mathrm{res}}$.

%------------------------------------------------------------
\subsection{Setup}
%------------------------------------------------------------

\begin{table}[h]
\centering
\begin{tabular}{L{4.2cm} L{7.2cm}}
\toprule
Models & \texttt{hc\_gcn\_l8}, \texttt{mhc\_gcn\_l8}, \texttt{mhc\_lite\_gcn\_l8} \\
Regimes & synthetic ID, synthetic OOD, Cora \\
Seeds & 5 \\
Depth analyzed & all 8 layers \\
Synthetic sample & 3 ID graphs + 3 OOD graphs per seed \\
Outputs & \texttt{outputs/analysis/hc\_matrices\_gcn\_l8/} \\
\bottomrule
\end{tabular}
\end{table}

%------------------------------------------------------------
\subsection{Matrix Statistics}
%------------------------------------------------------------

For each learned routing matrix $H_{\mathrm{res}} \in \mathbb{R}^{n_{\mathrm{str}}\times n_{\mathrm{str}}}$,
we report the following quantities.

\paragraph{Distance to identity.}
This measures how far the routing departs from ``no mixing'':
\[
d_I = \left\| H_{\mathrm{res}} - I \right\|_F
\]
Small values mean that each stream mostly keeps its own information.
Large values mean strong cross-stream remixing.

\paragraph{Row entropy.}
For row $i$, with weights $w_{ij}$:
\[
H_i = -\sum_j w_{ij}\log w_{ij}
\]
Low entropy means concentrated routing from one source stream.
High entropy means diffuse mixing from several source streams.

\paragraph{Row/column sum error.}
To test whether routing is approximately doubly stochastic, we measure
deviations from:
\[
\sum_j H_{ij} = 1,
\qquad
\sum_i H_{ij} = 1
\]
Small error means balanced transport; large error means unconstrained routing.

\paragraph{Permutation proximity.}
For double-stochastic routing, we also inspect proximity to permutation-like transport
using a Birkhoff-style decomposition:
\[
H_{\mathrm{res}} = \sum_k \lambda_k P_k,
\qquad
\lambda_k \ge 0,
\qquad
\sum_k \lambda_k = 1
\]
where $P_k$ are permutation matrices.
We track:
\begin{itemize}[noitemsep]
  \item \texttt{max\_permutation\_coeff}: dominant permutation weight
  \item \texttt{permutation\_coeff\_entropy}: concentration vs.\ diffusion of $\lambda_k$
\end{itemize}

%------------------------------------------------------------
\subsection{HC-GCN: progressively freer routing}
%------------------------------------------------------------

\begin{table}[h]
\centering
\caption{Representative HC-GCN routing statistics on synthetic OOD}
\begin{tabular}{C{1.5cm} R{3cm} R{3cm}}
\toprule
Layer & Identity distance & Row entropy \\
\midrule
0 & $\approx 0.064$ & $\approx 0.18$ \\
3 & $\approx 0.182$ & --- \\
7 & $\approx 0.951$ & $\approx 0.75$ \\
\bottomrule
\end{tabular}
\end{table}

\textbf{Observed behavior.}
\begin{itemize}[noitemsep]
  \item $H_{\mathrm{res}}$ moves steadily away from the identity as depth increases.
  \item Row entropy also grows with depth.
  \item Row/column sum errors become large in deeper layers.
\end{itemize}

\textbf{Interpretation.}
HC learns increasingly free, non-constrained inter-stream remixing.
The routing becomes both stronger and more diffuse with depth.
This is consistent with HC-GCN being the strongest variant on the synthetic benchmark:
the task benefits from aggressive multi-stream recombination of weak structural signals.

%------------------------------------------------------------
\subsection{mHC-GCN: strongly mixed but Sinkhorn-balanced}
%------------------------------------------------------------

\textbf{Observed behavior.}
\begin{itemize}[noitemsep]
  \item Before Sinkhorn, the raw residual matrix already wants to mix strongly:
  \[
  \texttt{raw\_identity\_distance} \approx 0.816 \quad \text{at layer 7}
  \]
  \item After Sinkhorn, the effective routing is much more controlled:
  \[
  \texttt{identity\_distance} \approx 0.403 \quad \text{at layer 7 on synthetic OOD}
  \]
  \item Row and column sum errors are essentially zero.
  \item This pattern is stable across synthetic ID, synthetic OOD, and Cora.
\end{itemize}

\textbf{Interpretation.}
mHC does \emph{not} learn weak mixing.
Instead, it learns strong mixing that is projected into a balanced transport plan.
The Sinkhorn constraint preserves more structure than HC and prevents unbalanced routing collapse.
This is consistent with mHC's stronger behavior on some real-graph transfer settings.

%------------------------------------------------------------
\subsection{mHC-lite-GCN: near-uniform default mixer}
%------------------------------------------------------------

\textbf{Typical values across layers and regimes.}
\begin{itemize}[noitemsep]
  \item $\texttt{identity\_distance} \approx 0.707$
  \item $\texttt{row\_abs\_entropy} \approx 1.0$
  \item $\texttt{max\_permutation\_coeff} \approx 0.50$
  \item $\texttt{permutation\_coeff\_entropy} \approx 1.0$
\end{itemize}

\textbf{Interpretation.}
With $n_{\mathrm{str}}=2$, mHC-lite behaves like a near-uniform convex combination
of identity and swap:
\[
H_{\mathrm{res}} \approx 0.5\,I + 0.5\,P_{\mathrm{swap}}
\]
The mechanism is highly regularized, weakly input-adaptive,
and nearly invariant with depth.
This explains why it is stable and competitive on large real graphs such as PubMed,
without being the strongest model on the synthetic benchmark.

%------------------------------------------------------------
\subsection{Mechanistic Summary}
%------------------------------------------------------------

\begin{table}[h]
\centering
\caption{Routing interpretation from learned $H_{\mathrm{res}}$ matrices}
\begin{tabular}{L{2.6cm} L{4.1cm} L{4.7cm}}
\toprule
Variant & Routing character & Consequence \\
\midrule
HC & Free, increasingly aggressive with depth &
Best synthetic OOD ranking; strongest remixing of structural signals \\
mHC & Strongly mixed, but Sinkhorn-balanced &
Better transfer stability; constrained transport prevents routing drift \\
mHC-lite & Near-uniform, almost depth-invariant mixer &
Stable, highly regularized behavior; good robustness on large real graphs \\
\bottomrule
\end{tabular}
\end{table}

\textbf{Main conclusion.}
The synthetic vs.\ real-graph ranking inversion is not a metric artefact.
It has a direct mechanistic signature in the routing matrices:
\begin{itemize}[noitemsep]
  \item synthetic inductive BC benefits from freer multi-stream routing,
  \item citation-graph transfer benefits more from constrained or highly regularized routing.
\end{itemize}

%------------------------------------------------------------
\section{Stage 9 --- EdgeConv Family Benchmark at L8 (2026-03-11)}
\label{sec:edgeconv}
%------------------------------------------------------------

\subsection{Motivation}

The GCN-family study shows that hyper-connections are useful on a simple backbone.
The next question is whether they still matter once the base convolution is already stronger.
The EdgeConv benchmark separates two effects:
\begin{enumerate}[noitemsep]
  \item backbone strength (\texttt{EdgeConv} vs.\ \texttt{GCN})
  \item routing mechanism (residual vs.\ HC/mHC/mHC-lite)
\end{enumerate}

%------------------------------------------------------------
\subsection{Setup}
%------------------------------------------------------------

\begin{table}[h]
\centering
\begin{tabular}{L{4.2cm} L{7.2cm}}
\toprule
Depth & $L=8$ \\
Seeds & 5 \\
New models & \texttt{edgeconv}, \texttt{edgeconv\_residual}, \texttt{hc\_edgeconv}, \texttt{mhc\_edgeconv}, \texttt{mhc\_lite\_edgeconv} \\
Training regime & matched to final GCN-family setup \\
Reference family & frozen GCN-family results at $L8$ \\
Output root & \texttt{outputs/followups/edgeconv\_vs\_gcn\_family\_l8/} \\
\bottomrule
\end{tabular}
\end{table}

%------------------------------------------------------------
\subsection{Main Results}
%------------------------------------------------------------

\begin{table}[h]
\centering
\caption{EdgeConv family benchmark at L8}
\begin{tabular}{L{4cm} R{2cm} R{2cm} R{2cm} R{2cm}}
\toprule
Model & OOD Kendall & OOD Spearman & Precision@10 & NDCG@10 \\
\midrule
\texttt{edgeconv\_l8}           & 0.3321 & 0.4698 & 0.2420 & 0.5276 \\
\texttt{edgeconv\_residual\_l8} & 0.4324 & 0.5955 & 0.2613 & 0.5718 \\
\texttt{hc\_edgeconv\_l8}       & 0.4292 & 0.5899 & 0.2467 & 0.5514 \\
\texttt{mhc\_edgeconv\_l8}      & \textbf{0.4360} & \textbf{0.5974} & 0.2553 & 0.5631 \\
\texttt{mhc\_lite\_edgeconv\_l8}& 0.4315 & 0.5927 & 0.2560 & 0.5621 \\
\bottomrule
\end{tabular}
\end{table}

For reference, the final GCN-family results at L8 are:

\begin{table}[h]
\centering
\caption{Reference GCN-family results at L8}
\begin{tabular}{L{4cm} R{2cm} R{2cm} R{2cm} R{2cm}}
\toprule
Model & OOD Kendall & OOD Spearman & Precision@10 & NDCG@10 \\
\midrule
\texttt{gcn\_l8}                & 0.3597 & 0.5055 & 0.2640 & 0.5632 \\
\texttt{gcn\_residual\_l8}      & 0.3685 & 0.5138 & 0.2673 & 0.5543 \\
\texttt{hc\_gcn\_l8}            & 0.4061 & 0.5576 & \textbf{0.2867} & 0.5669 \\
\texttt{mhc\_gcn\_l8}           & 0.3964 & 0.5479 & 0.2780 & 0.5678 \\
\texttt{mhc\_lite\_gcn\_l8}     & 0.3945 & 0.5442 & 0.2820 & 0.5610 \\
\bottomrule
\end{tabular}
\end{table}

%------------------------------------------------------------
\subsection{Interpretation}
%------------------------------------------------------------

\textbf{Observed behavior.}
\begin{itemize}[noitemsep]
  \item EdgeConv alone is not enough:
  \[
  \texttt{edgeconv\_l8} = 0.3321 < \texttt{gcn\_l8} = 0.3597
  \]
  on OOD Kendall.
  \item Adding a simple residual path changes the picture completely:
  \[
  \texttt{edgeconv\_residual\_l8} = 0.4324
  \]
  which is above every GCN-family model on global OOD ranking.
  \item Adding HC/mHC/mHC-lite on top of EdgeConv gives only small deltas:
  \[
  \texttt{mhc\_edgeconv\_l8} = 0.4360
  \quad \text{vs.} \quad
  \texttt{edgeconv\_residual\_l8} = 0.4324
  \]
  \item On top-K metrics, the best GCN variants remain stronger:
  \[
  \texttt{hc\_gcn\_l8} = 0.2867
  \quad \text{vs.} \quad
  \texttt{edgeconv\_residual\_l8} = 0.2613
  \]
  on Precision@10.
\end{itemize}

\textbf{Interpretation.}
The backbone matters more than initially hoped.
On GCN, hyper-connections are the main source of the gain.
On EdgeConv, the main source of the gain is the stronger convolution itself plus a standard residual path.
The marginal utility of HC/mHC/mHC-lite becomes much smaller once the backbone is already strong.

\textbf{Main conclusion.}
EdgeConv + residual is stronger than the final GCN family on global OOD ranking,
but the comparison is not like-for-like in terms of model class and capacity.
HC-GCN remains competitive if the target metric is top-K retrieval rather than global ranking.

%------------------------------------------------------------
\section{Consolidated Conclusions}
%------------------------------------------------------------

\begin{enumerate}
  \item \textbf{Hyper-connections improve inductive betweenness ranking on simple backbones.}
        On the synthetic multi-graph benchmark, the GCN family improves from
        plain GCN to HC/mHC/mHC-lite, and the gain is not reducible to a
        standard residual connection.

  \item \textbf{HC-GCN is the strongest variant on the main synthetic protocol.}
        At $L=8$, it achieves the best OOD Kendall within the GCN family:
        \[
          \text{HC-GCN}\ (0.406) > \text{mHC-GCN}\ (0.396)
          \approx \text{mHC-lite-GCN}\ (0.395) > \text{GCN}\ (0.360)
        \]
        The useful regime is moderate depth ($L4$--$L8$), not very deep scaling.

  \item \textbf{The internal routing mechanisms are genuinely different.}
        Matrix analysis shows that HC learns increasingly free and aggressive
        routing, mHC learns strong but Sinkhorn-balanced transport, and
        mHC-lite behaves like a highly regularized near-uniform mixer.
        These patterns are consistent with the observed performance differences.

  \item \textbf{The ranking between HC, mHC, and mHC-lite is distribution-dependent.}
        On the synthetic benchmark, HC is best. On real citation graphs,
        mHC is strongest on Cora and CiteSeer, while mHC-lite is strongest on
        PubMed. The family transfers well, but no single variant dominates
        every regime.

  \item \textbf{Backbone choice remains decisive.}
        EdgeConv with a standard residual path outperforms the final GCN family
        on global OOD ranking, while bringing only marginal additional gains
        from HC/mHC/mHC-lite. This suggests that hyper-connections are most
        useful as boosters for simpler backbones, not as universal replacements
        for stronger architectures.
\end{enumerate}

%------------------------------------------------------------
\section{Limits}
%------------------------------------------------------------

\begin{itemize}
  \item Very deep scaling ($L > 8$) is not validated on this task.
  \item Real-graph transfer is evaluated on a small number of citation graphs
        and should be read as external validation rather than as the main
        protocol.
  \item The routing analysis supports mechanistic interpretations, but remains
        an experimental interpretation rather than a formal proof.
  \item Capacity differs substantially across backbone families
        (e.g.\ GCN vs.\ EdgeConv / GCNII / APPNP), so cross-family comparisons
        must be interpreted with care.
\end{itemize}

\end{document}
